\documentclass[11pt,letterpaper,openany]{book}

% fonts and text
\usepackage[T1]{fontenc}
\usepackage[utf8]{inputenc}
\usepackage[sc]{mathpazo}
\renewcommand{\ttdefault}{txtt}
\linespread{1.06}
\usepackage{microtype}
\usepackage{amsmath,amssymb}

% page
\usepackage[left=1.15in,right=1.15in,top=1.05in,bottom=1.05in,headheight=15pt]{geometry}
\raggedbottom
\setlength{\emergencystretch}{2.5em}
\setlength{\columnsep}{24pt}

% colour
\usepackage{xcolor}
\definecolor{navy}{RGB}{24,52,96}
\definecolor{slate}{RGB}{96,106,124}
\definecolor{hairline}{RGB}{196,201,210}

% packages
\usepackage{tcolorbox}
\tcbuselibrary{skins,breakable}
\usepackage{titlesec}
\usepackage{fancyhdr}
\usepackage{multicol}
\usepackage{enumitem}
\usepackage{booktabs,tabularx,array}
\usepackage{float}
\usepackage{pifont}
\usepackage{tikz}
\usepackage{listings}
\lstset{basicstyle=\ttfamily\small,columns=fullflexible,keepspaces=true,showstringspaces=false}
\usepackage[hidelinks,bookmarksnumbered,pdftitle={Algebra 1: A Programmer's Introduction},pdfauthor={CS-Based Mathematics Curriculum}]{hyperref}

% headings
\titleformat{\chapter}[display]
  {\normalfont}
  {\color{slate}\large\scshape Chapter \thechapter}
  {8pt}
  {\color{navy}\fontsize{27}{31}\selectfont\bfseries}
  [\vspace{8pt}{\color{navy}\titlerule[1.1pt]}]
\titleformat{name=\chapter,numberless}[display]
  {\normalfont}
  {}
  {0pt}
  {\color{navy}\fontsize{27}{31}\selectfont\bfseries}
  [\vspace{8pt}{\color{navy}\titlerule[1.1pt]}]
\titlespacing*{\chapter}{0pt}{10pt}{34pt}
\titleformat{\section}{\normalfont\Large\bfseries\color{navy}}{\thesection}{0.8em}{}
\titlespacing*{\section}{0pt}{3.4ex plus 1ex minus .2ex}{1.4ex plus .2ex}
\titleformat{\subsection}{\normalfont\large\bfseries\color{navy}}{\thesubsection}{0.8em}{}
\titleformat{\paragraph}[runin]{\normalfont\bfseries\color{navy}}{}{0pt}{}[.]
\titlespacing*{\paragraph}{0pt}{1.6ex plus .6ex minus .2ex}{0.55em}

% running heads
\renewcommand{\chaptermark}[1]{\markboth{Chapter \thechapter\quad #1}{}}
\renewcommand{\sectionmark}[1]{\markright{\thesection\quad #1}}
\renewcommand{\headrule}{{\color{hairline}\hrule width\headwidth height\headrulewidth}}
\fancypagestyle{body}{\fancyhf{}
  \fancyhead[LE]{\small\thepage\qquad\textcolor{slate}{\leftmark}}
  \fancyhead[RO]{\small\textcolor{slate}{\rightmark}\qquad\thepage}
  \renewcommand{\headrulewidth}{0.3pt}}
\fancypagestyle{plain}{\fancyhf{}\fancyfoot[C]{\small\thepage}\renewcommand{\headrulewidth}{0pt}}

% boxes
\newtcolorbox{keybox}[1]{enhanced,breakable,sharp corners,
  colback=navy!5,colframe=navy!65,boxrule=0.5pt,
  title=#1,fonttitle=\bfseries\small,colbacktitle=navy!88!black,coltitle=white,
  left=10pt,right=10pt,top=7pt,bottom=7pt,before skip=12pt,after skip=14pt}
\newtcolorbox{examplebox}{enhanced,breakable,frame hidden,sharp corners,colback=white,
  borderline west={2.5pt}{0pt}{navy!55},left=10pt,right=2pt,top=3pt,bottom=3pt,
  before skip=10pt,after skip=10pt}
\newcounter{example}[chapter]
\renewcommand{\theexample}{\thechapter.\arabic{example}}
\newenvironment{example}{\refstepcounter{example}\begin{examplebox}\noindent\textcolor{navy}{\textbf{\textsc{Example~\theexample}}}\quad\ignorespaces}{\end{examplebox}}
\newcommand{\solution}{\par\smallskip\noindent\textit{Solution.}\enspace\ignorespaces}

% lists
\newlist{exercises}{enumerate}{1}
\setlist[exercises]{label=\textbf{\thechapter.\arabic*},leftmargin=*,labelsep=0.6em,itemsep=5pt,parsep=0pt,topsep=4pt}
\newlist{codeexercises}{enumerate}{1}
\setlist[codeexercises]{label=\textbf{C\thechapter.\arabic*},leftmargin=*,labelsep=0.6em,itemsep=7pt,parsep=0pt,topsep=4pt}
\newlist{reviewexercises}{enumerate}{1}
\setlist[reviewexercises]{label=\textbf{E.\arabic*},leftmargin=*,labelsep=0.6em,itemsep=5pt,parsep=0pt,topsep=4pt}
\newlist{checklist}{itemize}{1}
\setlist[checklist]{label=\ding{111},leftmargin=*,labelsep=0.7em,itemsep=4pt}
\setlist[itemize]{leftmargin=*,itemsep=3pt}
\setlist[enumerate]{leftmargin=*,itemsep=3pt}

% helpers
\newcommand{\chapterintro}[1]{\noindent{\large\itshape\color{slate}#1}\par\vspace{14pt}}
\newcommand{\answer}[2]{\par\noindent\hangindent=2.4em\hangafter=1\textbf{#1}\enspace #2}
\newcommand{\lerp}{\operatorname{lerp}}
\newcommand{\clamp}{\operatorname{clamp}}

\begin{document}
\frontmatter

% ================= TITLE PAGE =================
\begin{titlepage}
\centering
\vspace*{1.6in}
{\color{slate}\scshape\large CS-Based Mathematics Curriculum\par}
{\color{slate}\scshape\large Course 1\par}
\vspace{0.6in}
{\color{navy}\bfseries\fontsize{48}{54}\selectfont Algebra 1\par}
\vspace{0.3in}
{\color{navy}\rule{0.6\textwidth}{1.2pt}}\par
\vspace{0.3in}
{\Large\itshape A Programmer's Introduction\par}
\vspace{0.25in}
{\large With exercises to solve by hand and then in code\par}
\vfill
{\large Pilot Edition\par}
{\large October 2026\par}
\vspace{0.6in}
\end{titlepage}

% ================= FRONT MATTER =================
\pagestyle{plain}

\chapter*{Preface: How to Use This Book}
\markboth{Preface}{}
\addcontentsline{toc}{chapter}{Preface: How to Use This Book}

Every chapter runs the same way: read the core ideas, work the examples with the solutions covered, solve the exercises by hand, and only then write the code. The programming exercises reuse the hand exercises as test cases, so an unsolved exercise leaves you with nothing to test against.

Each chapter has four parts:

\begin{itemize}
\item \textbf{Core ideas}: the rules and definitions, with the computer science reason for each.
\item \textbf{Worked examples}: cover the solution, attempt it, then compare.
\item \textbf{Exercises}: numbered 1.1, 1.2 and so on; pencil, paper, no calculator, no code.
\item \textbf{Programming exercises}: numbered C1.1, C1.2 and so on; Python~3, plain functions, \lstinline!assert! statements as tests.
\end{itemize}

Rules that make the book work:

\begin{enumerate}
\item Do the exercises in one sitting before opening an editor. Check them against the answer key at the back of the book, and log every miss with the rule you broke.
\item Treat each programming exercise as a unit test of your understanding. Name the function exactly as given so later chapters can import it.
\item Keep all code in one folder, one file per chapter (\lstinline!ch01.py!, \lstinline!ch02.py!). Chapter~11 assembles them into the course project.
\item Use only the standard library. \lstinline!fractions!, \lstinline!math! and \lstinline!random! are allowed; NumPy and SymPy are not, because the point is to write what they do.
\item An exercise marked \textbf{(proof)} is answered in full sentences, not symbols alone.
\end{enumerate}

The book covers the ten topics of Course~1 in the curriculum, in the same order, and ends with the course project and an exit review. Plan on about one chapter per week.

\cleardoublepage
\tableofcontents

% ================= MAIN MATTER =================
\mainmatter
\pagestyle{body}

\chapter{Numbers and Operations}

\chapterintro{The real numbers are the values algebra works with, and the first job is knowing which kind you have and how the operations on them are ordered. Every one of these rules has a twin in your language's number types and operator precedence.}

\section{Core Ideas}

\paragraph{Number sets} Natural numbers ($1, 2, 3, \ldots$), integers ($\ldots, -2, -1, 0, 1, \ldots$), rationals (any fraction $p/q$ with integer $p$ and nonzero integer $q$), irrationals (cannot be written as a fraction: $\sqrt{2}$, $\pi$), and reals (rationals and irrationals together). A decimal that ends or repeats is rational; one that does neither is irrational. In code: integers are exact, floats are a finite approximation of reals with about 16 significant digits.

\paragraph{Order of operations} Parentheses, then exponents, then multiplication and division left to right, then addition and subtraction left to right. A tower of exponents is read from the top down: $2^{3^2}$ means $2^9 = 512$, not $(2^3)^2 = 64$. A negative sign without parentheses is applied after the exponent, so $-3^2 = -9$ while $(-3)^2 = 9$. Programming languages follow the same precedence; \lstinline!**! in Python is also right-associative.

\paragraph{Absolute value} $|x|$ is the distance from $x$ to $0$, so $|x| = x$ when $x$ is zero or positive and $|x| = -x$ when $x$ is negative. The distance between two numbers $a$ and $b$ is $|a - b|$.

\paragraph{Exponents and roots} $a^n$ is $n$ copies of $a$ multiplied together. $a^0 = 1$ for any nonzero $a$, and $a^{-n} = 1/a^n$. The square root of $a$ is the nonnegative number whose square is $a$.

\paragraph{Quotient and remainder} Dividing $a$ by $b$ gives a quotient $q$ and remainder $r$ with $a = bq + r$ and $0 \le r < b$. In Python, \lstinline!17 // 5! is 3 and \lstinline!17 % 5! is 2. Languages disagree on negative inputs: Python's \lstinline!-7 % 3! is 2 (floor division), JavaScript's is $-1$ (truncation).

\paragraph{Floating point} $0.1$ and $0.2$ have no exact binary representation, so \lstinline!0.1 + 0.2 == 0.3! is false. Compare floats with a tolerance or use exact fractions.

\section{Worked Examples}

\begin{example}
Evaluate $3 + 4 \times 2^2 - (6 - 1)$.
\solution Exponent first: $2^2 = 4$. Then $4 \times 4 = 16$. Parentheses: $6 - 1 = 5$. Then $3 + 16 - 5 = 14$.
\end{example}

\begin{example}
Simplify $|{-7}| - |3 - 10|$.
\solution $|{-7}| = 7$ and $|3 - 10| = |{-7}| = 7$, so the result is $0$.
\end{example}

\begin{example}
Write $0.375$ as a fraction.
\solution $0.375 = 375/1000$. The greatest common factor of $375$ and $1000$ is $125$, so $375/1000 = 3/8$.
\end{example}

\section{Exercises}

\begin{multicols}{2}
\begin{exercises}
\item Evaluate $18 - 2 \times 3^2 + 4 \div 2$.
\item Evaluate $(-3)^2$ and $-3^2$, and explain the difference in one sentence.
\item Evaluate $|4 - 9| + |{-2}| \times 3$.
\item Order from least to greatest: $-\tfrac{3}{4}$, $-0.8$, $-\tfrac{2}{3}$, $-0.7$.
\item Write $0.125$, $2.6$ and $0.\overline{3}$ as fractions in lowest terms.
\item Evaluate $2^{3^2}$ and $(2^3)^2$.
\item Find the quotient $q$ and remainder $r$ with $47 = 6q + r$ and $0 \le r < 6$.
\item Classify each as rational or irrational: $\sqrt{49}$, $\sqrt{50}$, $0.1010010001\ldots$, $22/7$.
\item Evaluate $5^0 + 2^{-2} + \sqrt{81}$.
\item A thermometer reads $-12$ degrees at dawn and $9$ degrees at noon. What is the change in temperature, and what is the distance between the two readings?
\end{exercises}
\end{multicols}

\section{Programming Exercises}

\begin{codeexercises}
\item Show in Python that \lstinline!0.1 + 0.2 == 0.3! is false. Then make the comparison true two ways: with \lstinline!math.isclose!, and with \lstinline!fractions.Fraction!. Write one sentence on why the float version fails.
\item Write \lstinline!floor_divmod(a, b)! that returns the quotient and remainder with $0 \le r < |b|$ for any integers $a$ and nonzero $b$, without using \lstinline!//!, \lstinline!%! or \lstinline!divmod!. Test it against \lstinline!divmod! for every $a$ from $-10$ to $10$ and $b$ in $\{3, -3\}$. Note which cases would differ in JavaScript.
\item Write \lstinline!abs_value(x)! without calling \lstinline!abs!, and \lstinline!distance(a, b)! for two numbers on a number line. Assert that \lstinline!distance(-12, 9)! matches your answer to Exercise~1.10.
\item Write \lstinline!gcd(a, b)! using the Euclidean algorithm, then \lstinline!to_fraction(s)! that turns a terminating decimal string such as \lstinline!"0.375"! into a reduced fraction as a tuple of integers. Check it against \lstinline!fractions.Fraction! for the values in Exercise~1.5.
\end{codeexercises}

\chapter{Expressions and Simplification}

\chapterintro{An algebraic expression is a formula with variables in it, and simplifying one is the same act as refactoring code: the result must evaluate to the same value for every input. The tools are combining like terms and the distributive property, and the mental model is the expression tree.}

\section{Core Ideas}

\paragraph{Terms and coefficients} An expression is a sum of terms. In $3x^2 - 5x + 7$, the terms are $3x^2$, $-5x$ and $7$; the coefficients are $3$, $-5$ and $7$ (the constant term). Like terms have the same variable raised to the same power: $3x$ and $-5x$ are like terms, $3x$ and $3x^2$ are not.

\paragraph{Combining like terms} Add the coefficients, keep the variable part: $3x + 5x = 8x$, and $4x^2 - x^2 = 3x^2$. This is just the distributive property run backward: $3x + 5x = (3 + 5)x$.

\paragraph{Distributive property} $a(b + c) = ab + ac$. It works with subtraction and negatives: $-(x - 4) = -x + 4$, and $2(3x - 5) = 6x - 10$. Run backward it is factoring out a common factor: $6x - 10 = 2(3x - 5)$.

\paragraph{Evaluating} Replace each variable with its value, wrap negatives in parentheses, then follow the order of operations.

\paragraph{Expression trees} Every expression is a tree: operators at the branches, numbers and variables at the leaves. $2 \times (x + 3)$ has $\times$ at the root with children $2$ and $(x + 3)$, as in Figure~\ref{fig:tree}. The tree makes precedence explicit, which is why parsers build one before evaluating.

\begin{figure}[H]
\centering
\begin{tikzpicture}[level distance=11mm,level 1/.style={sibling distance=22mm},level 2/.style={sibling distance=14mm},
  every node/.style={circle,draw=navy!60,fill=navy!5,minimum size=7mm,inner sep=1pt,font=\small},
  edge from parent/.style={draw=navy!60}]
\node{$\times$}
  child{node{$2$}}
  child{node{$+$}
    child{node{$x$}}
    child{node{$3$}}};
\end{tikzpicture}
\caption{The expression tree for $2 \times (x + 3)$. The parentheses are not stored; the shape carries them.}
\label{fig:tree}
\end{figure}

\paragraph{Translating words} ``5 less than twice a number'' is $2n - 5$, not $5 - 2n$. ``The quotient of $a$ and $b$'' is $a/b$. Write the variable's meaning down before writing the expression.

\section{Worked Examples}

\begin{example}
Simplify $3(2x - 4) + 5x - 7$.
\solution Distribute: $6x - 12 + 5x - 7$. Combine like terms: $11x - 19$.
\end{example}

\begin{example}
Evaluate $2a^2 - 3ab$ at $a = -2$ and $b = 3$.
\solution Substitute: $2(-2)^2 - 3(-2)(3) = 2(4) + 18 = 26$.
\end{example}

\begin{example}
Simplify $-(x - 4) + 2(3 - x)$.
\solution Distribute both: $-x + 4 + 6 - 2x$. Combine: $-3x + 10$.
\end{example}

\section{Exercises}

\begin{multicols}{2}
\begin{exercises}
\item Simplify $4x + 7 - 2x + 3$.
\item Simplify $5(2y - 3) - (y + 4)$.
\item Simplify $3x^2 + 2x - x^2 + 5x - 8$.
\item Evaluate $x^2 - 4x + 1$ at $x = -3$.
\item Evaluate $\dfrac{a + b}{a - b}$ at $a = 5$ and $b = -1$.
\item A rectangle has width $w$ and a length 3 more than twice its width. Write a simplified expression for its perimeter.
\item Factor out the greatest common factor of $12x^2 + 18x$.
\item Simplify $2(x + 3y) - 3(2x - y) + 4y$.
\item Draw the expression tree for $(x + 2) \times (x - 5)$ and for $x + 2 \times x - 5$. Evaluate both at $x = 4$.
\item A phone plan costs \$20 per month plus \$0.05 per text and \$0.10 per minute. Write the monthly cost for $t$ texts and $m$ minutes, then find the cost for 300 texts and 120 minutes.
\end{exercises}
\end{multicols}

\section{Programming Exercises}

\begin{codeexercises}
\item Represent a polynomial in one variable as a dictionary from power to coefficient, so $3x^2 + 2x - 8$ is \lstinline!{2: 3, 1: 2, 0: -8}!. Write \lstinline!combine(terms)! that takes a list of \lstinline!(coefficient, power)! pairs and returns the combined dictionary with zero coefficients removed. Test it with Exercise~2.3.
\item Write \lstinline!evaluate(poly, x)! for that dictionary form. Assert it agrees with your hand answer to Exercise~2.4.
\item Build an expression tree with classes \lstinline!Num!, \lstinline!Var!, \lstinline!Add!, \lstinline!Sub! and \lstinline!Mul!, each with an \lstinline!eval(env)! method where \lstinline!env! maps variable names to values. Encode both trees from Exercise~2.9 and confirm the two values you computed by hand.
\item Write \lstinline!to_string(poly)! that prints the dictionary form in standard order with correct signs, so \lstinline!{2: 2, 1: 7, 0: -8}! prints as \lstinline!2x^2 + 7x - 8! and \lstinline!{1: 1, 0: -1}! prints as \lstinline!x - 1!. Test every edge case you can think of: coefficient 1, coefficient $-1$, missing powers, a constant alone.
\end{codeexercises}

\chapter{Linear Equations in One Variable}

\chapterintro{Solving an equation means undoing, in reverse order, everything that was done to the variable, while keeping both sides equal. Checking the answer by substitution is the first unit test you will ever write by hand.}

\section{Core Ideas}

\paragraph{The balance rule} An equation stays true if you do the same thing to both sides: add, subtract, multiply or divide by the same nonzero number. Inverse operations undo each other: subtraction undoes addition, division undoes multiplication.

\paragraph{A solving procedure} Simplify each side (distribute, combine like terms). Move variable terms to one side and constants to the other. Divide by the coefficient of the variable. Check by substituting into the original equation.

\paragraph{Three outcomes} Most equations have one solution. If the variables cancel and you are left with a false statement such as $0 = 7$, there is no solution. If you are left with a true statement such as $0 = 0$, every number is a solution (an identity). Your solver in Exercise~C3.1 must return all three.

\paragraph{Literal equations} Solving $P = 2l + 2w$ for $w$ uses the same moves and gives $w = (P - 2l)/2$. This is how you turn one formula into its inverse, which you will do constantly when converting units or reversing a transformation.

\paragraph{Proportions} $a/b = c/d$ is equivalent to $ad = bc$ (cross-multiplication), provided $b$ and $d$ are nonzero.

\paragraph{Word problems} Name the unknown with its unit, write the relationship as an equation, solve, then check the answer against the story, not just the equation.

\section{Worked Examples}

\begin{example}
Solve $3x - 7 = 11$.
\solution Add 7: $3x = 18$. Divide by 3: $x = 6$. Check: $3(6) - 7 = 11$.
\end{example}

\begin{example}
Solve $2(x + 3) = 5x - 9$.
\solution Distribute: $2x + 6 = 5x - 9$. Subtract $2x$: $6 = 3x - 9$. Add 9: $15 = 3x$. So $x = 5$. Check: $2(8) = 16$ and $25 - 9 = 16$.
\end{example}

\begin{example}
Solve $A = \tfrac{1}{2}bh$ for $h$.
\solution Multiply by 2: $2A = bh$. Divide by $b$: $h = 2A/b$.
\end{example}

\section{Exercises}

\begin{multicols}{2}
\begin{exercises}
\item Solve $5x + 4 = 29$.
\item Solve $\dfrac{x}{3} - 2 = 4$.
\item Solve $4(x - 2) = 2x + 10$.
\item Solve $7 - 2(x + 1) = 3x - 5$.
\item Solve $3(2x - 1) = 6x + 4$. State what kind of outcome this is.
\item Solve $2(x + 4) - x = x + 8$. State what kind of outcome this is.
\item Solve $F = \tfrac{9}{5}C + 32$ for $C$.
\item Solve $\dfrac{3}{x} = \dfrac{12}{20}$.
\item Solve $\dfrac{x + 1}{4} = \dfrac{2x - 3}{2}$. Give the exact answer as a fraction.
\item A consultant charges a \$150 setup fee plus \$85 per hour. A client was billed \$830. How many hours were worked?
\item Two trains start 300 miles apart and travel toward each other, one at 60 mph and the other at 40 mph. After how many hours do they meet?
\end{exercises}
\end{multicols}

\section{Programming Exercises}

\begin{codeexercises}
\item Write \lstinline!solve_linear(a, b, c, d)! that solves $ax + b = cx + d$ and returns one of: \lstinline!("one", x)!, \lstinline!("none",)! or \lstinline!("all",)!. Test it against Exercises~3.1 through 3.6.
\item Write \lstinline!check_linear(a, b, c, d, x)! that substitutes $x$ into both sides and compares them with \lstinline!math.isclose!. Use it to verify every \lstinline!("one", x)! result from C3.1, so your solver and your checker test each other.
\item Write \lstinline!c_to_f(c)! and \lstinline!f_to_c(f)! from Exercise~3.7. Assert that \lstinline!f_to_c(c_to_f(c))! returns \lstinline!c! (within tolerance) for every integer $c$ from $-50$ to $150$.
\item Rewrite C3.1 to use \lstinline!fractions.Fraction! for all arithmetic so that Exercise~3.9 returns exactly $7/3$ rather than $2.3333$. Decide what the function should do if the inputs are floats.
\end{codeexercises}

\chapter{Linear Functions and Graphs}

\chapterintro{A linear function changes at a constant rate, and that rate, the slope, is the whole story: it is the pixels-per-unit of a scale, the speed of a moving sprite, and the step of a linear interpolation. This chapter turns lines into functions you can graph, build from two points, and use to map one range of values onto another.}

\section{Core Ideas}

\paragraph{Slope} For two points $(x_1, y_1)$ and $(x_2, y_2)$, the slope is the change in $y$ over the change in $x$.

\begin{keybox}{Slope of a line}
\[
m = \frac{y_2 - y_1}{x_2 - x_1}
\]
\end{keybox}

\paragraph{Three forms of a line} Slope-intercept $y = mx + b$ ($b$ is the $y$-intercept). Point-slope $y - y_1 = m(x - x_1)$, the fastest way to write a line from a point and a slope. Standard $Ax + By = C$ with integer coefficients. Convert between them with the moves from Chapter~3.

\paragraph{Graphing} Plot the $y$-intercept, then use the slope as rise over run to find a second point. For standard form, find both intercepts by setting $x = 0$ and then $y = 0$.

\paragraph{Parallel and perpendicular} Parallel lines have equal slopes. Perpendicular lines have slopes whose product is $-1$, so the perpendicular slope is the negative reciprocal.

\paragraph{Linear interpolation} The value a fraction $t$ of the way from $a$ to $b$ is $\lerp(a, b, t) = a + (b - a)t$. At $t = 0$ it is $a$, at $t = 1$ it is $b$. Mapping a value from one range to another is a line through two points.

\begin{keybox}{Mapping one range onto another}
\[
\operatorname{map}(x) = \text{out}_{\min} + (x - \text{in}_{\min})\,\frac{\text{out}_{\max} - \text{out}_{\min}}{\text{in}_{\max} - \text{in}_{\min}}
\]
\end{keybox}

\paragraph{Direct variation} $y = kx$ is a line through the origin; $k$ is the constant of variation and also the slope.

\section{Worked Examples}

\begin{example}
Find the line through $(2, 5)$ and $(6, 13)$.
\solution Slope: $(13 - 5)/(6 - 2) = 2$. Point-slope: $y - 5 = 2(x - 2)$. Simplified: $y = 2x + 1$.
\end{example}

\begin{example}
Graph $3x + 2y = 12$.
\solution Set $y = 0$: $x = 4$, so $(4, 0)$. Set $x = 0$: $y = 6$, so $(0, 6)$. Draw the line through both.
\end{example}

\begin{example}
Compute $\lerp(10, 20, 0.25)$.
\solution $10 + (20 - 10)(0.25) = 12.5$.
\end{example}

\section{Exercises}

\begin{multicols}{2}
\begin{exercises}
\item Find the slope of the line through $(-1, 4)$ and $(3, -4)$.
\item Write the equation of the line with slope $\tfrac{3}{2}$ through $(4, -1)$ in slope-intercept form.
\item Write the equation of the line through $(0, 3)$ and $(5, 0)$.
\item Find the $x$- and $y$-intercepts of $4x - 5y = 20$.
\item Write $y = -2x + 7$ in standard form with integer coefficients.
\item Write the line parallel to $y = 4x - 1$ through $(1, 1)$.
\item Write the line perpendicular to $y = 4x - 1$ through $(1, 1)$.
\item A chart maps a data value $v$ in the range $0$ to $100$ onto a pixel position $x$ in the range $40$ to $440$. Write $x$ as a linear function of $v$, and find the pixel for $v = 35$.
\item Celsius to Fahrenheit is linear and passes through $(0, 32)$ and $(100, 212)$. Derive the formula from these two points.
\item Compute $\lerp(-5, 15, 0.3)$. Then find the $t$ for which $\lerp(-5, 15, t) = 10$.
\item A taxi charges \$3.50 plus \$2.25 per mile. Write the cost as a function of distance, find the cost of a 7.2-mile ride, and find the distance a \$30 fare buys (to one decimal).
\end{exercises}
\end{multicols}

\section{Programming Exercises}

\begin{codeexercises}
\item Write \lstinline!lerp(a, b, t)! and \lstinline!inverse_lerp(a, b, v)!, which returns the $t$ for which lerp gives $v$. Assert they round-trip for a grid of values, and check Exercise~4.10.
\item Write \lstinline!map_range(x, in_min, in_max, out_min, out_max)! using the two functions from C4.1, and verify Exercise~4.8 with it.
\item Write \lstinline!line_through(p1, p2)! returning \lstinline!(m, b)! for a non-vertical line, or \lstinline!("vertical", x)! when the two points share an $x$-coordinate. Test it with the points in Exercises~4.1 through 4.3 and with a vertical pair.
\item Plot $y = \tfrac{3}{2}x - 7$ and your perpendicular line from Exercise~4.7 on the same axes with matplotlib (the one allowed exception to the standard-library rule) or an HTML canvas. Observe that they only look perpendicular when both axes use the same scale, and explain why.
\item Write \lstinline!ease_in_quad(t)! $= t \times t$ and print a table comparing it to \lstinline!lerp(0, 100, t)! for $t$ in $0, 0.25, 0.5, 0.75, 1$. Which one is a linear function, and how can you tell from the table alone?
\end{codeexercises}

\chapter{Systems of Linear Equations}

\chapterintro{A system of two linear equations asks where two lines meet, and the answer is a point, no point, or every point on a shared line. The elimination method you learn here is the $2 \times 2$ case of Gaussian elimination, which is the workhorse of Course~6.}

\section{Core Ideas}

\paragraph{What a solution is} A pair $(x, y)$ that satisfies both equations at once. Graphically it is the intersection of the two lines.

\paragraph{Substitution} Solve one equation for a variable, substitute that expression into the other equation, solve, then back-substitute. Best when one equation is already solved for a variable.

\paragraph{Elimination} Multiply one or both equations by constants so that one variable has opposite coefficients, then add the equations to eliminate it. These are row operations: scaling a row and adding one row to another never changes the solution set.

\paragraph{Three outcomes} Lines with different slopes cross once (one solution). Parallel distinct lines never cross (no solution; elimination produces a false statement like $0 = 5$). The same line written twice gives infinitely many solutions (elimination produces $0 = 0$).

\paragraph{The determinant test} For $a_1 x + b_1 y = c_1$ and $a_2 x + b_2 y = c_2$, the number $D = a_1 b_2 - a_2 b_1$ decides it. If $D$ is not zero there is exactly one solution, given by Cramer's rule.

\begin{keybox}{Cramer's rule for a $2 \times 2$ system}
\[
x = \frac{c_1 b_2 - c_2 b_1}{a_1 b_2 - a_2 b_1}, \qquad y = \frac{a_1 c_2 - a_2 c_1}{a_1 b_2 - a_2 b_1}
\]
\end{keybox}

\paragraph{Back substitution} A triangular system such as $x + y + z = 6$, $2y + z = 7$, $z = 3$ is solved from the bottom up. Elimination's whole purpose is to produce this shape.

\section{Worked Examples}

\begin{example}
Solve $y = 2x - 1$ and $3x + y = 9$ by substitution.
\solution Replace $y$: $3x + (2x - 1) = 9$, so $5x = 10$ and $x = 2$. Then $y = 2(2) - 1 = 3$. Check: $6 + 3 = 9$.
\end{example}

\begin{example}
Solve $2x + 3y = 12$ and $4x - 3y = 6$ by elimination.
\solution The $y$-coefficients are already opposites, so add: $6x = 18$, $x = 3$. Substitute: $6 + 3y = 12$, $y = 2$.
\end{example}

\begin{example}
Where does the line $y = x$ cross $y = -x + 6$?
\solution Set $x = -x + 6$, so $2x = 6$ and $x = 3$; the point is $(3, 3)$. This is exactly the computation a 2D engine does to find where two paths cross.
\end{example}

\section{Exercises}

\begin{multicols}{2}
\begin{exercises}
\item Solve by substitution: $y = x + 4$ and $2x + y = 10$.
\item Solve by elimination: $x + y = 7$ and $x - y = 1$.
\item Solve by elimination: $3x + 2y = 16$ and $2x - 5y = -2$. Show the multipliers you chose.
\item Classify the system $2x - 4y = 8$ and $x - 2y = 4$ as one solution, none, or infinitely many, and explain using slopes.
\item Classify $3x + y = 5$ and $6x + 2y = 7$ the same way.
\item Find the intersection of $y = 2x - 3$ and $y = -x + 6$.
\item Twelve tickets cost \$105. Adult tickets are \$10 and child tickets are \$7. How many of each were sold?
\item You need 20 liters of a 30\% solution by mixing a 10\% solution with a 50\% solution. How many liters of each?
\item For the system in Exercise~5.3, compute $D$ and solve it again with Cramer's rule. Confirm the answers match.
\item Solve by back substitution: $x + y + z = 6$, $2y + z = 7$, $z = 3$.
\end{exercises}
\end{multicols}

\section{Programming Exercises}

\begin{codeexercises}
\item Write \lstinline!solve_2x2(a1, b1, c1, a2, b2, c2)! using Cramer's rule, returning \lstinline!("one", x, y)!, \lstinline!("none",)! or \lstinline!("infinite",)!. When $D$ is zero, decide between none and infinite by checking whether the two equations are proportional. Test with Exercises~5.1 through 5.6.
\item Write \lstinline!intersect(p1, p2, p3, p4)! that takes two lines, each given by two points, and returns their intersection point or \lstinline!None! when they are parallel. Build each line's standard-form coefficients from its two points and reuse C5.1. Verify Exercise~5.6 by picking two points on each line.
\item Write \lstinline!eliminate(rows)! that takes the augmented matrix \lstinline![[a1, b1, c1], [a2, b2, c2]]!, performs elimination with explicit row operations, and prints the matrix after each operation. Run it on Exercise~5.3 and compare the printed steps with your hand work.
\item Write a brute-force checker that scans all integer pairs $(x, y)$ from $-20$ to $20$ and reports every pair satisfying both equations. Use it to confirm C5.1 on every exercise with integer answers, and to show that Exercise~5.4 has many solutions while Exercise~5.5 has none in that range.
\end{codeexercises}

\chapter{Inequalities}

\chapterintro{An inequality describes a range rather than a single value, and ranges are what every bounds check, clamp and validation rule in your code enforces. The solving moves are the same as for equations, with one exception that catches everyone at least once.}

\section{Core Ideas}

\paragraph{The one new rule} Adding or subtracting on both sides keeps the direction. Multiplying or dividing both sides by a negative number reverses it: $-2x < 6$ becomes $x > -3$. Dividing by a positive number keeps the direction.

\paragraph{Writing the answer} On a number line, an open circle means the endpoint is excluded ($<$ or $>$) and a closed circle means it is included ($\le$ or $\ge$). Interval notation uses a round bracket for excluded and a square bracket for included: $x \ge 2$ is $[2, \infty)$, and $2 \le x < 7$ is $[2, 7)$.

\paragraph{Compound inequalities} ``And'' means both conditions hold, which is the overlap of two ranges; $-2 \le 2x - 6 < 8$ is solved on all three parts at once. ``Or'' means at least one holds, which is the union of two ranges.

\paragraph{Absolute value inequalities} $|x| < a$ means $-a < x < a$ ($x$ is within $a$ of zero). $|x| > a$ means $x < -a$ or $x > a$ ($x$ is at least $a$ away from zero). Rewrite, then solve the compound inequality. If $a$ is negative, $|x| < a$ has no solution and $|x| > a$ is true for every $x$.

\paragraph{Clamping} $\clamp(x, \text{lo}, \text{hi}) = \max(\text{lo}, \min(x, \text{hi}))$ returns the nearest value satisfying $\text{lo} \le x \le \text{hi}$. A range check and a clamp are the two things code does with an inequality: test it or enforce it.

\section{Worked Examples}

\begin{example}
Solve $-3x + 4 > 13$.
\solution Subtract 4: $-3x > 9$. Divide by $-3$ and flip: $x < -3$.
\end{example}

\begin{example}
Solve $-2 \le 2x - 6 < 8$.
\solution Add 6 everywhere: $4 \le 2x < 14$. Divide by 2: $2 \le x < 7$, or $[2, 7)$.
\end{example}

\begin{example}
Solve $|x - 5| \le 3$.
\solution Rewrite: $-3 \le x - 5 \le 3$. Add 5: $2 \le x \le 8$. Every $x$ within 3 of 5.
\end{example}

\section{Exercises}

\begin{multicols}{2}
\begin{exercises}
\item Solve $4x - 9 \ge 7$.
\item Solve $5 - 2x < 11$.
\item Solve $3(x + 2) \le 5x - 4$.
\item Solve $-1 < \dfrac{x + 3}{2} \le 4$ and write the answer in interval notation.
\item Solve $x - 2 < -6$ or $2x \ge 10$.
\item Solve $|2x + 1| < 7$.
\item Solve $|x - 4| \ge 2$.
\item Solve $|3x| + 2 < 1$.
\item Write in interval notation: (a) $x \ge -2$ and $x < 3$; (b) $x < 0$ or $x > 9$.
\item A ride admits riders between 42 and 54 inches tall, inclusive. Write this as a compound inequality and as a single absolute value inequality.
\item An API costs \$0.004 per call after a \$10 monthly fee, and the monthly budget is \$50. What is the largest number of calls you can make?
\end{exercises}
\end{multicols}

\section{Programming Exercises}

\begin{codeexercises}
\item Write \lstinline!clamp(x, lo, hi)!. For 10{,}000 random floats, assert that $\text{lo} \le \clamp(x, \text{lo}, \text{hi}) \le \text{hi}$ and that clamp returns $x$ unchanged whenever $x$ was already in range.
\item Write \lstinline!solve_inequality(a, b, op, c)! for $ax + b \mathrel{\text{op}} c$, where \lstinline!op! is one of \lstinline!"<"!, \lstinline!"<="!, \lstinline!">"!, \lstinline!">="!. Return the solution as an interval tuple \lstinline!(lo, hi, lo_closed, hi_closed)! using \lstinline!float("-inf")! and \lstinline!float("inf")! for open ends. It must flip the operator when $a$ is negative. Test with Exercises~6.1 through 6.3.
\item Write \lstinline!contains(interval, x)! for that interval tuple. Test the two intervals from Exercise~6.9 at their endpoints and just inside and outside each.
\item Write \lstinline!solve_abs_inequality(a, b, op, c)! for $|ax + b| \mathrel{\text{op}} c$, returning one interval, a list of two intervals, or an empty list. Test with Exercises~6.6 through 6.8.
\item Write \lstinline!can_ride(h)! using the absolute value form from Exercise~6.10, and a second function \lstinline!can_ride_compound(h)! that uses the compound form. Assert they agree for every integer height from 30 to 60.
\end{codeexercises}

\chapter{Exponents and Exponential Growth}

\chapterintro{Exponents are repeated multiplication, and the difference between adding a constant each step and multiplying by one is the difference between a loop that finishes and one that does not. This chapter also covers powers of 2, which is how memory, bits and the growth of algorithms are measured.}

\section{Core Ideas}

\paragraph{The exponent rules} For nonzero $a$ and $b$ and integers $m$ and $n$, the six rules below follow directly from exponents being repeated multiplication.

\begin{keybox}{Rules of exponents}
\[
a^m a^n = a^{m+n} \qquad \frac{a^m}{a^n} = a^{m-n} \qquad (a^m)^n = a^{mn}
\]
\[
(ab)^n = a^n b^n \qquad a^0 = 1 \qquad a^{-n} = \frac{1}{a^n}
\]
\end{keybox}

\paragraph{Scientific notation} A number written as $c \times 10^n$ with $1 \le c < 10$. In code it is \lstinline!6.02e23!. Multiply the leading numbers and add the exponents, then renormalize.

\paragraph{Exponential growth and decay} A quantity that grows by a fixed percentage $r$ each period follows $y = a(1 + r)^t$, where $a$ is the starting amount. Decay is $y = a(1 - r)^t$. Doubling every period is $y = a \cdot 2^t$.

\paragraph{Powers of 2} $2^{10} = 1024$, which is why a kilobyte is ``about a thousand'' bytes. $2^{20}$ is about a million, $2^{30}$ about a billion, $2^{32} = 4{,}294{,}967{,}296$ (the number of values in a 32-bit integer), and $2^{64}$ is about $1.8 \times 10^{19}$.

\paragraph{Binary} Each binary digit is a power of 2: $101101$ means $32 + 8 + 4 + 1 = 45$. To represent $n$ distinct values you need the smallest $k$ with $2^k \ge n$ bits.

\paragraph{Linear versus exponential} Adding a fixed amount each step gives a line; multiplying by a fixed factor gives a curve that eventually outruns any line. An algorithm that takes $2^n$ steps is unusable past $n$ of about 40, however fast the machine.

\section{Worked Examples}

\begin{example}
Simplify $(2x^3 y)^2 \cdot x^{-4}$.
\solution Square each factor: $4x^6 y^2$. Multiply by $x^{-4}$: $4x^2 y^2$.
\end{example}

\begin{example}
\$1000 grows at 5\% per year for 3 years. What is it worth?
\solution $1000 \times 1.05^3 = 1000 \times 1.157625 = \$1157.63$.
\end{example}

\begin{example}
Convert $101101$ (binary) to decimal.
\solution Place values from the right are $1, 2, 4, 8, 16, 32$. The ones are at 32, 8, 4 and 1, so the value is $45$.
\end{example}

\section{Exercises}

\begin{multicols}{2}
\begin{exercises}
\item Simplify $x^5 \cdot x^{-2} \cdot x$.
\item Simplify $(3a^2 b^{-1})^3$ with positive exponents only.
\item Simplify $\dfrac{12x^7}{4x^9}$.
\item Evaluate $2^{-3} \cdot 4^2$.
\item Write $0.00042$ and $3{,}500{,}000$ in scientific notation, then multiply them and give the product in scientific notation.
\item A culture of 500 bacteria doubles every 3 hours. How many are there after 24 hours?
\item A \$2000 laptop loses 20\% of its value each year. What is it worth after 4 years?
\item Write $2^{10}$, $2^{16}$ and $2^{32}$ in decimal and state the nearest power of 10 for each.
\item How many bits are needed to represent 1000 distinct values? 70{,}000 distinct values?
\item Convert $10011$ (binary) to decimal, and $77$ (decimal) to binary.
\item One algorithm takes $n$ steps and another takes $2^n$ steps. Compute both for $n = 10$, $20$ and $30$. What is the smallest $n$ for which $2^n$ exceeds one million?
\end{exercises}
\end{multicols}

\section{Programming Exercises}

\begin{codeexercises}
\item Write \lstinline!power(base, n)! for any integer $n$, negative included, with a simple loop. Then write \lstinline!fast_power(base, n)! using squaring (if $n$ is even, square the result for $n/2$; if odd, multiply by the base once more). Count the multiplications each version makes for $n = 1000$.
\item Write \lstinline!bits_needed(n)! that returns the smallest $k$ with $2^k \ge n$. Verify it against \lstinline!(n - 1).bit_length()! for every $n$ from 2 to 100{,}000 and explain in a comment why that expression works.
\item Write \lstinline!to_binary(n)! by repeated division by 2 and \lstinline!from_binary(s)! by place values, without \lstinline!bin()! or \lstinline!int(s, 2)!. Assert they round-trip for every $n$ from 0 to 5000, and check Exercise~7.10.
\item Starting from \$1000, print a year-by-year table of a linear account that adds \$50 per year and an exponential account that grows 5\% per year. Find the first year the exponential account is ahead, and the first year it is ahead by more than \$1000.
\item Compute \lstinline!2**53 + 1! as an integer and as a float. Show that the float version equals \lstinline!2**53!, and explain the result in one sentence using the 53-bit precision of a double.
\end{codeexercises}

\chapter{Polynomials and Factoring}

\chapterintro{A polynomial is a sum of terms with whole-number powers, and factoring is the reverse of multiplying it out. Multiplication is mechanical; factoring is a search, which is why it makes a good first algorithm to write.}

\section{Core Ideas}

\paragraph{Vocabulary} The degree is the highest power. The leading coefficient is the coefficient of that term. Standard form lists terms from highest power to lowest. A monomial has one term, a binomial two, a trinomial three.

\paragraph{Adding and subtracting} Combine like terms. For subtraction, distribute the minus sign across every term of the second polynomial first.

\paragraph{Multiplying} Every term of the first polynomial multiplies every term of the second, then combine like terms. For two binomials this is FOIL (first, outer, inner, last). Three products are worth memorizing.

\begin{keybox}{Special products}
\[
(a + b)^2 = a^2 + 2ab + b^2 \qquad (a - b)^2 = a^2 - 2ab + b^2 \qquad (a + b)(a - b) = a^2 - b^2
\]
\end{keybox}

\paragraph{Factoring, in order of attempts} First pull out the greatest common factor. Then for $x^2 + bx + c$, find two numbers that multiply to $c$ and add to $b$. For $ax^2 + bx + c$, find two numbers that multiply to $ac$ and add to $b$, split the middle term, and factor by grouping. Check for a difference of squares or a perfect square trinomial. Four terms usually factor by grouping.

\paragraph{Coefficient arrays} A polynomial is a list of coefficients indexed by power: $2x^3 - 6x^2 + 2x - 1$ is \lstinline![-1, 2, -6, 2]! with index 0 holding the constant. Addition is elementwise; multiplication of degree $m$ and degree $n$ gives degree $m + n$, where the coefficient at index $k$ is the sum of all products of coefficients whose indices add to $k$.

\paragraph{Horner's method} Evaluate from the leading coefficient inward: $2x^3 - 6x^2 + 2x - 1 = ((2x - 6)x + 2)x - 1$. One multiplication and one addition per coefficient, with no exponentiation at all.

\section{Worked Examples}

\begin{example}
Multiply $(2x - 3)(x + 5)$.
\solution First $2x^2$, outer $10x$, inner $-3x$, last $-15$. Combined: $2x^2 + 7x - 15$.
\end{example}

\begin{example}
Factor $x^2 - 5x - 14$.
\solution Two numbers with product $-14$ and sum $-5$ are $-7$ and $2$. So $(x - 7)(x + 2)$.
\end{example}

\begin{example}
Factor $6x^2 + 7x - 3$.
\solution $ac = -18$; two numbers with product $-18$ and sum $7$ are $9$ and $-2$. Split: $6x^2 + 9x - 2x - 3$. Group: $3x(2x + 3) - 1(2x + 3) = (3x - 1)(2x + 3)$.
\end{example}

\section{Exercises}

\begin{multicols}{2}
\begin{exercises}
\item Add $(3x^2 - 2x + 5) + (x^2 + 4x - 9)$.
\item Subtract $(5x^3 - x + 2) - (2x^3 + 3x^2 - 4)$.
\item Multiply $(x + 4)(x - 6)$.
\item Expand $(3x - 2)^2$.
\item Multiply $(2x + 5)(2x - 5)$.
\item Multiply $(x - 1)(x^2 + x + 1)$.
\item Factor $8x^3 - 12x^2$.
\item Factor $x^2 + 9x + 20$.
\item Factor $x^2 - 3x - 28$.
\item Factor $2x^2 + 11x + 12$.
\item Factor $49y^2 - 64$.
\item Factor $x^3 + 2x^2 - 9x - 18$ completely.
\item Evaluate $2x^3 - 6x^2 + 2x - 1$ at $x = 3$ with Horner's method, writing down each intermediate value.
\end{exercises}
\end{multicols}

\section{Programming Exercises}

\begin{codeexercises}
\item Using coefficient lists, write \lstinline!poly_add(p, q)!, \lstinline!poly_sub(p, q)! and \lstinline!poly_mul(p, q)!. Strip trailing zeros so the degree is honest. Test with Exercises~8.1 through 8.6.
\item Write \lstinline!horner(coeffs, x)! and \lstinline!naive_eval(coeffs, x)! (the version that calls \lstinline!x ** k! for each term). Count the multiplications each makes for a degree-50 polynomial, treating \lstinline!x ** k! as $k - 1$ multiplications.
\item Write \lstinline!factor_trinomial(a, b, c)! for integer coefficients. Search every pair of integer factors of $ac$ for a pair that sums to $b$, then factor by grouping, and return the two binomials as coefficient lists or \lstinline!None! if there is no integer factorization. Test with Exercises~8.8 through 8.10.
\item Write \lstinline!poly_to_string(coeffs)! that prints standard form, so \lstinline![-24, -2, 1]! prints as \lstinline!x^2 - 2x - 24!. Reuse or improve your C2.4.
\item For every factoring exercise in this chapter, multiply your factors back together with \lstinline!poly_mul! and assert the result equals the original. A factorization that does not survive this test is wrong, whatever the answer key says.
\end{codeexercises}

\chapter{Quadratics}

\chapterintro{A quadratic is the simplest curve, and solving one tells you when a thrown object lands, when two moving things collide, and where a parabola peaks. There are four solving methods, and the last one, the quadratic formula, is the other three automated.}

\section{Core Ideas}

\paragraph{The shape} $y = ax^2 + bx + c$ graphs as a parabola, opening up when $a > 0$ and down when $a < 0$. The $y$-intercept is $c$. The vertex (the highest or lowest point) is at $x = -b/(2a)$, and the vertical line through it is the axis of symmetry.

\paragraph{Zero product property} If a product is zero, at least one factor is zero. So if $(x - 3)(x - 4) = 0$, then $x = 3$ or $x = 4$. This is why factoring solves equations: set one side to zero, factor, set each factor to zero.

\paragraph{Square roots} If $x^2 = k$ with $k > 0$, then $x$ is the positive or negative square root of $k$. The same works for $(x - 4)^2 = 25$.

\paragraph{Completing the square} For $x^2 + bx$, adding $(b/2)^2$ makes a perfect square $(x + b/2)^2$. Rewrite the equation that way, then take square roots. This also converts standard form to vertex form $y = a(x - h)^2 + k$, where $(h, k)$ is the vertex.

\paragraph{The quadratic formula} For $ax^2 + bx + c = 0$ with $a$ nonzero, the solutions are given directly by the formula below; Exercise~9.12 asks you to derive it.

\begin{keybox}{The quadratic formula}
\[
x = \frac{-b \pm \sqrt{b^2 - 4ac}}{2a}
\]
\end{keybox}

\paragraph{The discriminant} The quantity $b^2 - 4ac$ under the root decides how many real solutions there are: positive means two, zero means one (a repeated root), negative means none. Your solver should report this before it tries to take a square root.

\paragraph{Motion} An object launched upward follows $h(t) = -16t^2 + v_0 t + h_0$ in feet (or $-4.9t^2 + v_0 t + h_0$ in meters). Its peak is at the vertex and it lands when $h = 0$.

\section{Worked Examples}

\begin{example}
Solve $x^2 - 7x + 12 = 0$.
\solution Factor: $(x - 3)(x - 4) = 0$. So $x = 3$ or $x = 4$.
\end{example}

\begin{example}
Solve $x^2 + 6x - 7 = 0$ by completing the square.
\solution Move the constant: $x^2 + 6x = 7$. Add $(6/2)^2 = 9$ to both sides: $x^2 + 6x + 9 = 16$. So $(x + 3)^2 = 16$, $x + 3 = \pm 4$, and $x = 1$ or $x = -7$.
\end{example}

\begin{example}
Solve $2x^2 + 3x - 4 = 0$ by the formula.
\solution Discriminant: $9 + 32 = 41$. $x = \dfrac{-3 \pm \sqrt{41}}{4}$, which is about $0.851$ and $-2.351$.
\end{example}

\section{Exercises}

\begin{multicols}{2}
\begin{exercises}
\item Solve $x^2 - 9 = 0$.
\item Solve $x^2 + 5x = 0$.
\item Solve $x^2 - 2x - 15 = 0$ by factoring.
\item Solve $(x - 4)^2 = 25$.
\item Solve $x^2 + 8x + 3 = 0$ by completing the square. Leave the answer exact.
\item Solve $3x^2 - 5x - 2 = 0$ by the quadratic formula.
\item Find the discriminant and the number of real solutions for each: $x^2 + 4x + 5 = 0$; $4x^2 - 12x + 9 = 0$; $2x^2 - x - 3 = 0$.
\item Find the vertex and axis of symmetry of $y = 2x^2 - 8x + 3$.
\item Write $y = x^2 - 6x + 11$ in vertex form.
\item A ball is thrown upward at 48 feet per second from a height of 6 feet, so $h(t) = -16t^2 + 48t + 6$. When is it highest, how high does it get, and when does it hit the ground (to two decimals)?
\item A rectangle has a perimeter of 40. Write its area as a function of its width, then find the dimensions that make the area largest.
\item \textbf{(proof)} Starting from $ax^2 + bx + c = 0$ with $a$ nonzero, derive the quadratic formula by completing the square. Every step must be justified in a sentence.
\end{exercises}
\end{multicols}

\section{Programming Exercises}

\begin{codeexercises}
\item Write \lstinline!solve_quadratic(a, b, c)! returning \lstinline!(count, roots)! where \lstinline!count! is 0, 1 or 2 and \lstinline!roots! is a sorted list. If $a$ is zero, delegate to \lstinline!solve_linear! from Chapter~3. Test with Exercises~9.1 through 9.7.
\item Run C9.1 on $a = 1$, $b = -100{,}000{,}000$, $c = 1$. One root is wrong. Implement the stable version: compute $q = -\bigl(b + \operatorname{sign}(b)\sqrt{D}\bigr)/2$, then $x_1 = q/a$ and $x_2 = c/q$. Compare both versions and explain the term ``catastrophic cancellation'' in a comment.
\item Write \lstinline!vertex(a, b, c)! and plot $y = 2x^2 - 8x + 3$ with its vertex marked. Confirm Exercise~9.8.
\item Simulate Exercise~9.10 in steps of $0.01$ seconds, advancing the height with the formula each step, and report the first time the height is at or below zero. Compare with the exact answer from the formula.
\item A sprite moves along $x(t) = 2t^2 + 3t$ and a wall stands at $x = 50$. Use C9.1 to find the collision time, then confirm it by stepping the simulation.
\end{codeexercises}

\chapter{Functions}

\chapterintro{A function assigns exactly one output to every input, which is the same promise a pure function makes in code, and it is the promise that makes caching, testing and composition safe. This chapter makes the idea precise and gives you composition and inverses, which the rest of the curriculum uses constantly.}

\section{Core Ideas}

\paragraph{Definition} A function is a rule that gives each input exactly one output. The domain is the set of allowed inputs; the range is the set of outputs actually produced. A relation that sends one input to two different outputs is not a function, which is what the vertical line test checks on a graph.

\paragraph{Notation} $f(x) = 3x - 2$ names the function $f$ and its rule. $f(4)$ means ``the output when the input is 4,'' which is $10$. $f(a + 1)$ means substitute $a + 1$ for every $x$.

\paragraph{Four representations} A table of input-output pairs, a graph, a formula, and a program. All four describe the same object; being able to move between them is the skill.

\paragraph{Domain restrictions} Division by zero and square roots of negatives are undefined in the reals, so $f(x) = 1/(x - 3)$ has domain all $x$ except $3$, and $g(x) = \sqrt{x - 4}$ has domain $x \ge 4$. In code these are the inputs that raise exceptions.

\paragraph{Composition} $(f \circ g)(x) = f(g(x))$: apply $g$ first, then $f$. Order matters, and $f(g(x))$ and $g(f(x))$ are usually different functions. Composition is what a pipeline of transformations is.

\paragraph{Inverses} The inverse $f^{-1}$ undoes $f$, so $f(f^{-1}(x)) = x$ and $f^{-1}(f(x)) = x$. To find it, swap $x$ and $y$ in the formula and solve for $y$. A function has an inverse only when no two inputs share an output.

\paragraph{Piecewise functions} Different rules on different parts of the domain. In code this is an if-else chain, and the conditions must not overlap.

\section{Worked Examples}

\begin{example}
Let $f(x) = 3x - 2$ and $g(x) = x^2$. Find $f(g(2))$ and $g(f(2))$.
\solution $f(g(2)) = f(4) = 10$, while $g(f(2)) = g(4) = 16$.
\end{example}

\begin{example}
Find the domain of $f(x) = 1/(x - 3)$.
\solution The denominator is zero when $x = 3$, so the domain is every real number except $3$.
\end{example}

\begin{example}
Shipping costs \$5 for the first pound and \$2 for each additional pound. Write the cost function and find the cost of a 3.5-pound package.
\solution $c(w) = 5$ if $w \le 1$, and $c(w) = 5 + 2(w - 1)$ if $w > 1$. For a 3.5-pound package, $c(3.5) = 5 + 2(2.5) = 10$.
\end{example}

\section{Exercises}

\begin{multicols}{2}
\begin{exercises}
\item For $f(x) = 2x^2 - x + 1$, find $f(0)$, $f(-2)$ and $f(a + 1)$, simplified.
\item Which of these are functions: (a) $\{(1, 2)$, $(2, 3)$, $(1, 4)\}$; (b) $\{(1, 2)$, $(2, 2)$, $(3, 2)\}$; (c) $y^2 = x$?
\item Find the domain of $g(x) = \sqrt{x - 4}$, and of $h(x) = \dfrac{x + 1}{x^2 - 9}$.
\item For $f(x) = x + 5$ and $g(x) = 2x$, find $f(g(x))$ and $g(f(x))$.
\item For $f(x) = x^2 - 1$ and $g(x) = 3x$, find $(f \circ g)(2)$ and $(g \circ f)(2)$.
\item Find the inverse of $f(x) = 4x - 3$ and verify that $f(f^{-1}(x)) = x$.
\item Let $f(x) = x^2$ for $x < 0$ and $f(x) = 2x + 1$ for $x \ge 0$. Find $f(-3)$, $f(0)$ and $f(4)$.
\item Find the range of $f(x) = x^2 + 2$ and of $g(x) = -|x| + 3$.
\item Temperature in Fahrenheit is $F(C) = 1.8C + 32$, and a comfort score is $s(F) = -|F - 72| + 100$. Write $s$ as a function of $C$ and evaluate it at 20 degrees Celsius.
\item \textbf{(proof)} Explain why $f(x) = x^2$ has no inverse on all real numbers, and why it does have one when the domain is restricted to $x \ge 0$.
\end{exercises}
\end{multicols}

\section{Programming Exercises}

\begin{codeexercises}
\item Write \lstinline!compose(f, g)! that returns a new function computing $f(g(x))$. Test Exercises~10.4 and 10.5 with lambdas, and show that \lstinline!compose(f, g)! and \lstinline!compose(g, f)! differ.
\item Write \lstinline!with_domain(f, allowed)! that returns a function raising \lstinline!ValueError! when \lstinline!allowed(x)! is false and otherwise returns $f(x)$. Wrap the two functions from Exercise~10.3 and test an input inside and outside each domain.
\item Represent a piecewise function as a list of \lstinline!(condition, rule)! pairs and write \lstinline!piecewise(rules)! that returns a function applying the first matching rule. Encode Exercise~10.7 and confirm your three values. Add a check that raises if no rule matches.
\item Write \lstinline!is_inverse(f, f_inv, samples)! that returns true when $f(f_{\text{inv}}(x))$ and $f_{\text{inv}}(f(x))$ both equal $x$ within tolerance for every sample. Test Exercise~10.6, and show that it fails for $f(x) = x^2$ with $f_{\text{inv}}(x) = \sqrt{x}$ when the samples include negatives.
\item Write \lstinline!memoize(f)! that caches results in a dictionary. Explain in a comment why this is only correct because $f$ is a function in the mathematical sense, and give an example of a Python callable it would break.
\end{codeexercises}

\chapter{Course Project}

\chapterintro{The project is a small algebra engine that reads expressions as text, simplifies and evaluates them, solves linear equations, $2 \times 2$ systems and quadratics, and plots the results. Almost every piece already exists in your chapter files; this chapter assembles them behind a parser and a command line.}

Build it in six stages, each with its own acceptance tests. Do not start a stage until the previous one passes. Because the engine reads and prints plain text, exponents in its input and output use the caret, as in \lstinline!3x^2 + 2x - 5!.

\section{The Six Stages}

\begin{table}[H]
\centering
\footnotesize
\renewcommand{\arraystretch}{1.25}
\begin{tabularx}{\textwidth}{@{}>{\raggedright\arraybackslash}p{0.85in} >{\raggedright\arraybackslash}X >{\raggedright\arraybackslash}p{0.9in} >{\raggedright\arraybackslash}X@{}}
\toprule
\textbf{Stage} & \textbf{Build} & \textbf{Reuses} & \textbf{Done when} \\
\midrule
1. Tokenizer & \texttt{tokenize(s)} turns \texttt{"3x\textasciicircum{}2 + 2x - 5"} into a list of tokens: numbers, the variable, operators, parentheses & Nothing new & Whitespace is ignored, \texttt{-5} and \texttt{- 5} both work, and bad characters raise an error \\
2. Parser & \texttt{parse(tokens)} builds an expression tree honoring precedence and parentheses & C2.3 tree classes & \texttt{2 * (x + 3)} and \texttt{2 * x + 3} produce different trees, and \texttt{2\textasciicircum{}3\textasciicircum{}2} groups right to left \\
3. Polynomial form & \texttt{to\_poly(tree)} flattens a tree with $+$, $-$, $*$ and integer powers of $x$ into a coefficient list; \texttt{to\_string} prints it & C2.1, C8.1, C8.4 & \texttt{(x + 2) * (x - 5)} prints as \texttt{x\textasciicircum{}2 - 3x - 10}, and \texttt{evaluate} agrees with the tree's own \texttt{eval} for ten random $x$ \\
4. Linear solver & \texttt{solve("3x - 7 = 11")} parses both sides, moves everything to one side, and reports the solution, none, or all & C3.1, C3.4 & Exercises 3.1 through 3.6 all return the correct outcome, and 3.9 returns $7/3$ exactly \\
5. System and quadratic solvers & \texttt{solve\_system("3x + 2y = 16", "2x - 5y = -2")} and \texttt{solve("x\textasciicircum{}2 - 2x - 15 = 0")} & C5.1, C9.1, C9.2 & Exercises 5.1 through 5.6 and 9.1 through 9.7 all pass, including the no-solution and infinite cases, and the discriminant is reported \\
6. Plotter & \texttt{plot("x\textasciicircum{}2 - 2x - 15", -5, 7)} draws the curve with its roots and vertex marked & C4.4, C9.3 & A linear expression draws as a line and a quadratic as a parabola with the correct intercepts \\
\bottomrule
\end{tabularx}
\caption{The six stages of the algebra engine, what each reuses, and its acceptance test.}
\end{table}

\section{Acceptance Checklist}

\begin{checklist}
\item Every exercise in Chapters 2, 3, 5, 8 and 9 that has a numeric answer is reproduced by the engine from its text form.
\item A test file runs all of those checks with one command and prints a pass count.
\item The parser rejects malformed input with a clear message rather than a stack trace.
\item Arithmetic uses \lstinline!Fraction! throughout, with floats appearing only when a square root is taken.
\item A short README explains how to run it and which chapter each module came from.
\end{checklist}

\section{Stretch Goals}

Support a second variable in Stage~3 so systems parse from text. Add \lstinline!simplify! so \lstinline!2(x + 3) - x! prints as \lstinline!x + 6!. Print the step-by-step solution for a linear equation the way you would write it by hand.


% ================= BACK MATTER =================
\backmatter
\pagestyle{body}
\chapter*{Exit Review}
\markboth{Exit Review}{}
\addcontentsline{toc}{chapter}{Exit Review}

Solve these twelve problems by hand in one sitting, with no notes, then verify every numeric answer with your own engine from Chapter~11. Eleven correct passes the course; fewer means rereading the chapters you missed and retaking it a week later. These answers are deliberately not in the key.

\begin{multicols}{2}
\begin{reviewexercises}
\item Evaluate $-2^2 + |5 - 11| \div 3 \times 2^3$.
\item Simplify $3(2x - 1) - 2(x - 4) + x$.
\item Solve $5(x - 2) = 3x + 6$.
\item Write the equation of the line through $(2, -3)$ and $(-4, 9)$.
\item Solve the system $4x - y = 11$ and $2x + 3y = 9$.
\item Solve $|3x - 6| \le 9$ and write the answer in interval notation.
\item Simplify $\dfrac{(2x^2 y^{-3})^2}{4xy}$ with positive exponents only.
\item Factor $3x^3 - 12x$ completely.
\item Solve $2x^2 - 7x + 3 = 0$.
\item For $f(x) = x^2 + 1$ and $g(x) = x - 2$, find $f(g(x))$, $g(f(x))$ and $(f \circ g)(5)$.
\item A rocket's height is $h(t) = -16t^2 + 96t$. At what times is it 128 feet high, and what is its maximum height?
\item A file holds $3{,}000{,}000$ records, each needing a unique integer ID. How many bits per ID are required? If each record also stores a 4-byte timestamp, how many megabytes (one million bytes each) do the timestamps alone take?
\end{reviewexercises}
\end{multicols}

\medskip
The curriculum's exit check for Course~1 also applies, and the review above covers it:

\begin{checklist}
\item Solve any linear equation or inequality in one variable and verify the answer by substitution.
\item Solve a $2 \times 2$ linear system three ways and explain when it has no solution or infinitely many.
\item Factor and solve quadratics, and derive the quadratic formula by completing the square.
\item Explain in one sentence why $2^{10}$ is close to $10^3$ and why that matters for memory sizes.
\end{checklist}

\chapter*{Answers to Exercises}
\markboth{Answers to Exercises}{}
\addcontentsline{toc}{chapter}{Answers to Exercises}

Check only after finishing a chapter's full set. A wrong answer goes in your mistake log with the rule you broke.

\medskip
\begin{multicols}{2}
\small
\setlength{\parskip}{2.5pt}

\subsection*{Chapter 1}
\answer{1.1}{$2$}
\answer{1.2}{$9$ and $-9$; the exponent is applied before the negative sign unless the parentheses include the sign.}
\answer{1.3}{$11$}
\answer{1.4}{$-0.8$, $-\tfrac{3}{4}$, $-0.7$, $-\tfrac{2}{3}$}
\answer{1.5}{$\tfrac{1}{8}$, $\tfrac{13}{5}$, $\tfrac{1}{3}$}
\answer{1.6}{$512$ and $64$}
\answer{1.7}{$q = 7$, $r = 5$}
\answer{1.8}{Rational (it equals 7), irrational, irrational, rational}
\answer{1.9}{$10.25$, or $\tfrac{41}{4}$}
\answer{1.10}{Change of $+21$ degrees; distance $21$}

\subsection*{Chapter 2}
\answer{2.1}{$2x + 10$}
\answer{2.2}{$9y - 19$}
\answer{2.3}{$2x^2 + 7x - 8$}
\answer{2.4}{$22$}
\answer{2.5}{$\tfrac{2}{3}$}
\answer{2.6}{$6w + 6$}
\answer{2.7}{$6x(2x + 3)$}
\answer{2.8}{$-4x + 13y$}
\answer{2.9}{$-6$ and $7$}
\answer{2.10}{$20 + 0.05t + 0.10m$; \$47}

\subsection*{Chapter 3}
\answer{3.1}{$x = 5$}
\answer{3.2}{$x = 18$}
\answer{3.3}{$x = 9$}
\answer{3.4}{$x = 2$}
\answer{3.5}{No solution (the equation reduces to $-3 = 4$)}
\answer{3.6}{Identity; every real number is a solution}
\answer{3.7}{$C = \tfrac{5}{9}(F - 32)$}
\answer{3.8}{$x = 5$}
\answer{3.9}{$x = \tfrac{7}{3}$}
\answer{3.10}{8 hours}
\answer{3.11}{3 hours}

\subsection*{Chapter 4}
\answer{4.1}{$-2$}
\answer{4.2}{$y = \tfrac{3}{2}x - 7$}
\answer{4.3}{$y = -\tfrac{3}{5}x + 3$}
\answer{4.4}{$(5, 0)$ and $(0, -4)$}
\answer{4.5}{$2x + y = 7$}
\answer{4.6}{$y = 4x - 3$}
\answer{4.7}{$y = -\tfrac{1}{4}x + \tfrac{5}{4}$}
\answer{4.8}{$x = 4v + 40$; pixel $180$}
\answer{4.9}{$F = 1.8C + 32$}
\answer{4.10}{$1$; $t = 0.75$}
\answer{4.11}{$C(d) = 2.25d + 3.50$; \$19.70; about 11.8 miles}

\subsection*{Chapter 5}
\answer{5.1}{$(2, 6)$}
\answer{5.2}{$(4, 3)$}
\answer{5.3}{$(4, 2)$; for example multiply the first equation by 5 and the second by 2 to eliminate $y$}
\answer{5.4}{Infinitely many; both are the same line with slope $\tfrac{1}{2}$ and intercept $-2$}
\answer{5.5}{No solution; both have slope $-3$ with different intercepts, so they are parallel}
\answer{5.6}{$(3, 3)$}
\answer{5.7}{7 adult tickets and 5 child tickets}
\answer{5.8}{10 liters of each}
\answer{5.9}{$D = -19$; $x = 4$, $y = 2$}
\answer{5.10}{$x = 1$, $y = 2$, $z = 3$}

\subsection*{Chapter 6}
\answer{6.1}{$x \ge 4$}
\answer{6.2}{$x > -3$}
\answer{6.3}{$x \ge 5$}
\answer{6.4}{$-5 < x \le 5$, written $(-5, 5]$}
\answer{6.5}{$x < -4$ or $x \ge 5$}
\answer{6.6}{$-4 < x < 3$}
\answer{6.7}{$x \le 2$ or $x \ge 6$}
\answer{6.8}{No solution; an absolute value cannot be less than $-1$}
\answer{6.9}{(a) $[-2, 3)$; (b) $(-\infty, 0) \cup (9, \infty)$}
\answer{6.10}{$42 \le h \le 54$; $|h - 48| \le 6$}
\answer{6.11}{10{,}000 calls}

\subsection*{Chapter 7}
\answer{7.1}{$x^4$}
\answer{7.2}{$27a^6/b^3$}
\answer{7.3}{$3/x^2$}
\answer{7.4}{$2$}
\answer{7.5}{$4.2 \times 10^{-4}$ and $3.5 \times 10^{6}$; product $1.47 \times 10^{3}$}
\answer{7.6}{$128{,}000$}
\answer{7.7}{\$819.20}
\answer{7.8}{$1024$ (about $10^3$); $65{,}536$ (about $6.5 \times 10^4$); $4{,}294{,}967{,}296$ (about $4.3 \times 10^9$)}
\answer{7.9}{10 bits; 17 bits}
\answer{7.10}{$19$; $1001101$}
\answer{7.11}{$10$ vs $1024$; $20$ vs $1{,}048{,}576$; $30$ vs $1{,}073{,}741{,}824$; $n = 20$}

\subsection*{Chapter 8}
\answer{8.1}{$4x^2 + 2x - 4$}
\answer{8.2}{$3x^3 - 3x^2 - x + 6$}
\answer{8.3}{$x^2 - 2x - 24$}
\answer{8.4}{$9x^2 - 12x + 4$}
\answer{8.5}{$4x^2 - 25$}
\answer{8.6}{$x^3 - 1$}
\answer{8.7}{$4x^2(2x - 3)$}
\answer{8.8}{$(x + 4)(x + 5)$}
\answer{8.9}{$(x - 7)(x + 4)$}
\answer{8.10}{$(2x + 3)(x + 4)$}
\answer{8.11}{$(7y - 8)(7y + 8)$}
\answer{8.12}{$(x + 2)(x - 3)(x + 3)$}
\answer{8.13}{Intermediates $2, 0, 2, 5$; the value is $5$}

\subsection*{Chapter 9}
\answer{9.1}{$x = 3$ or $x = -3$}
\answer{9.2}{$x = 0$ or $x = -5$}
\answer{9.3}{$x = 5$ or $x = -3$}
\answer{9.4}{$x = 9$ or $x = -1$}
\answer{9.5}{$x = -4 \pm \sqrt{13}$}
\answer{9.6}{$x = 2$ or $x = -\tfrac{1}{3}$}
\answer{9.7}{$-4$ (none); $0$ (one); $25$ (two)}
\answer{9.8}{Vertex $(2, -5)$; axis $x = 2$}
\answer{9.9}{$y = (x - 3)^2 + 2$}
\answer{9.10}{Highest at $t = 1.5$ s, at 42 ft; hits the ground at about $t = 3.12$ s}
\answer{9.11}{$A(w) = w(20 - w)$; a 10 by 10 square, area 100}
\answer{9.12}{Divide by $a$; move $c/a$ to the right; add $\bigl(\tfrac{b}{2a}\bigr)^2$ to both sides; the left is now $\bigl(x + \tfrac{b}{2a}\bigr)^2$ and the right is $\tfrac{b^2 - 4ac}{4a^2}$; take square roots of both sides; subtract $\tfrac{b}{2a}$}

\subsection*{Chapter 10}
\answer{10.1}{$1$; $11$; $2a^2 + 3a + 2$}
\answer{10.2}{(a) no, 1 maps to two outputs; (b) yes; (c) no, $x = 4$ gives $y = 2$ and $y = -2$}
\answer{10.3}{$x \ge 4$; every $x$ except $3$ and $-3$}
\answer{10.4}{$2x + 5$; $2x + 10$}
\answer{10.5}{$35$; $9$}
\answer{10.6}{$f^{-1}(x) = \tfrac{x + 3}{4}$; then $4 \cdot \tfrac{x + 3}{4} - 3 = x$}
\answer{10.7}{$9$; $1$; $9$}
\answer{10.8}{$y \ge 2$; $y \le 3$}
\answer{10.9}{$s(C) = -|1.8C - 40| + 100$; $s(20) = 96$}
\answer{10.10}{$f(2) = f(-2) = 4$, so an inverse would have to send 4 to two inputs; restricted to $x \ge 0$ each output comes from exactly one input and the inverse is the square root}

\end{multicols}


\end{document}
